\section{Recovery of the coefficients}\label{sec:coefficient-recovery}

The holomorphic comparison from Proposition~\ref{prop:transport-matching}
is constrained by the geometry of its polarization bundle.  We first show
that this comparison is the identity, and then recover the coefficients from
the transport matrices.

\subsection{The bundle on the null conic}

Let $H\to\mathcal C$ be the holomorphic bundle with fiber $H_\theta$, and
let $\mathcal N\subset H$ be the null line subbundle with fiber
$\C\theta$.  Thus $\Ecal=H\oplus\mathcal O$, where $\mathcal O$ denotes
the trivial holomorphic line bundle on $\mathcal C$.
On $\mathbb P^1(\C)$, let $\mathcal O(-1)$ denote the tautological line
bundle and $\mathcal O(-2)$ its tensor square.

\begin{lemma}\label{lem:conic-bundle}
The bundle $H$ has no nonzero global holomorphic sections, and
$H/\mathcal N$ is holomorphically trivial.  More precisely, under the
identification $\mathcal C\simeq\mathbb P^1(\C)$,
\[
 H\simeq\mathcal O(-1)^{\oplus2},\qquad
 \mathcal N\simeq\mathcal O(-2),
\]
and their quotient sequence is
\begin{equation}\label{eq:conic-quotient}
 0\longrightarrow\mathcal O(-2)
 \xrightarrow{\ (z_0,z_1)^T\ }
 \mathcal O(-1)^{\oplus2}
 \xrightarrow{\ (-z_1,z_0)\ }
 \mathcal O\longrightarrow0.
\end{equation}
Here $[z_0:z_1]$ are homogeneous coordinates on $\mathbb P^1(\C)$.
\end{lemma}

\begin{proof}
A homogeneous parametrization of the conic is
\begin{equation}\label{eq:conic-parametrization}
 \theta(z_0,z_1)
   =(z_0^2-z_1^2,\ i(z_0^2+z_1^2),\ 2z_0z_1).
\end{equation}
It identifies $\mathbb P^1(\C)$ with $\mathcal C$: on the first coordinate
chart its inverse is
$z_1/z_0=\theta_3/(\theta_1-i\theta_2)$, and on the second it is
$z_0/z_1=\theta_3/(-\theta_1-i\theta_2)$.  The denominators cannot both
vanish at a point of $\mathcal C$.

Consider the two degree-one columns
\[
 h_1=(z_0,iz_0,z_1),\qquad
 h_2=(-z_1,iz_1,z_0).
\]
Direct calculation, with the complex bilinear product, gives
\begin{equation}\label{eq:conic-frame}
 \theta\cdot h_1=\theta\cdot h_2=0,\qquad
 h_1\times h_2=i\theta,\qquad
 \theta=z_0h_1+z_1h_2.
\end{equation}
To interpret these homogeneous formulas as bundle maps, write
$\mathcal O(-1)_{[z]}=\C z$ for $z=(z_0,z_1)$.
The assignment
\[
 (\xi_1z,\xi_2z)\longmapsto \xi_1h_1(z)+\xi_2h_2(z)
\]
is independent of the representative: replacing $z$ by $cz$ replaces
$\xi_j$ by $c^{-1}\xi_j$ and $h_j(z)$ by $ch_j(z)$.
Its local polynomial formulas are holomorphic. Since the images lie in
$H$ and are linearly independent at every point, it is an isomorphism
$\mathcal O(-1)^{\oplus2}\to H$. Likewise,
$\xi z^{\otimes2}\mapsto\xi\theta(z)$ is well defined because
$\theta(cz)=c^2\theta(z)$, and identifies $\mathcal O(-2)$ with
$\mathcal N$. The final identity in \eqref{eq:conic-frame} gives the
first map in \eqref{eq:conic-quotient}.

The expression $-z_1\xi_1+z_0\xi_2$ is also unchanged under the same
representative scaling, so its local polynomial formulas define a
holomorphic map to the trivial line bundle $\mathcal O$. The row
$(-z_1,z_0)$ is everywhere surjective, and its kernel is the image of
$(z_0,z_1)^T$.
This proves the quotient assertion, including at both coordinate-chart
endpoints.  It is also consistent with the induced bilinear form: for
$p=\xi_1 h_1+\xi_2 h_2$ in a local frame,
\[
 p\cdot p=(-z_1\xi_1+z_0\xi_2)^2.
\]

Finally, any holomorphic section of $H$ has holomorphic ambient components
on the compact conic, so those components are constant.  Such a constant
vector is orthogonal to every null vector.  The vectors
$\theta(1,0)$, $\theta(0,1)$ and $\theta(1,1)$ span $\C^3$, so the
section is zero.
\end{proof}

\begin{remark}
The exact sequence \eqref{eq:conic-quotient} does not split holomorphically:
a splitting would give a nonzero global section of $H$.  On the other
hand, the displayed splitting of $H$ itself gives nonscalar holomorphic
endomorphisms of $H$.  The next argument uses the transport equation to
restrict those endomorphisms.
\end{remark}

\begin{proposition}\label{prop:comparison-rigidity}
The comparison endomorphism of
Proposition~\ref{prop:transport-matching} is $G=\Id$.  Consequently, for
every $x\in\R^3$ and every nonzero null vector $\theta$,
\begin{equation}\label{eq:equal-restricted-transports}
 M_{1,\theta}=M_{2,\theta}\quad\text{on }\Ecal_\theta.
\end{equation}
\end{proposition}

\begin{proof}
Fix $x$.  Relative to $\Ecal=H\oplus\mathcal O$, the upper-right block
of the holomorphic endomorphism $G(x,\cdot)$ is a global holomorphic
section of $H$, and therefore vanishes by Lemma~\ref{lem:conic-bundle}.
Its lower-right block is a holomorphic scalar function on the compact
conic, and hence is independent of direction.  We may thus write
\[
 G=\begin{pmatrix}F&0\\ A&f\end{pmatrix},\qquad f=f(x).
\]
All these blocks are smooth in $x$.  Taking the upper-right block of
\eqref{eq:comparison-transport}, and using
\eqref{eq:transport-matrix}, gives
\[
 0=\theta f-F\theta.
\]
It follows that $F-f\Id_H$ vanishes on $\mathcal N$.  This holomorphic
bundle morphism consequently factors through $H/\mathcal N$.  By
Lemma~\ref{lem:conic-bundle}, its factor is a morphism
$\mathcal O\to H$, which must vanish.  Thus $F=f\Id_H$.

The upper-left block of \eqref{eq:comparison-transport} now reads
\begin{equation}\label{eq:comparison-rank}
 (D_\theta f)\Id_{H_\theta}=\theta A.
\end{equation}
The right side has rank at most one, while $H_\theta$ has dimension two.
Therefore $D_\theta f=0$.  Equation~\eqref{eq:comparison-rank} then gives
$A=0$, since $\theta\ne0$.  As null vectors span $\C^3$, we obtain
$\nabla f=0$.  The exterior identity value of $G$ fixes $f=1$, so $G=\Id$.
Substitution in \eqref{eq:comparison-transport} proves
\eqref{eq:equal-restricted-transports}.
\end{proof}

\subsection{A finite-type uniqueness lemma}

We will show that equality of the restricted bottom-left transport blocks
leaves at most a scalar-identity ambiguity, then use the following
elementary lemma to remove it. Indices
in the lemma range from $1$ to $n$; all sums are written explicitly.

\begin{lemma}\label{lem:finite-type}
Let $U\subset\R^n$ be connected and open, with $n\ge2$, and let
$\eta,q$ and $B_{ij}^{\ a}$ be smooth, real- or complex-valued functions
on $U$ satisfying
\begin{equation}\label{eq:finite-type-hypothesis}
 \partial_i\partial_j\eta
    =\sum_{a=1}^n B_{ij}^{\ a}\partial_a\eta+\delta_{ij}q.
\end{equation}
Then $(\nabla\eta,q)$ satisfies a homogeneous first-order system with
smooth coefficients.  In particular, if $\eta$ vanishes on a nonempty
open subset of $U$, then $\eta=q=0$ on $U$.
\end{lemma}

\begin{proof}
Put $v_a=\partial_a\eta$.  For each $s$, choose one index $i=i(s)$ with
$i\ne s$, without summing over $i$.  Equation~\eqref{eq:finite-type-hypothesis}
and commutation of third derivatives give
\begin{align*}
 \partial_s q
 &=\partial_i\left(\sum_{a=1}^n B_{si}^{\ a}v_a\right)
       -\partial_s\left(\sum_{a=1}^n B_{ii}^{\ a}v_a\right)\\
 &=\sum_{a=1}^n
       (\partial_i B_{si}^{\ a}-\partial_s B_{ii}^{\ a})v_a
    +\sum_{a=1}^n
       (B_{si}^{\ a}\partial_i v_a-B_{ii}^{\ a}\partial_s v_a).
\end{align*}
There is no derivative of $q$ on the right because the off-diagonal
entry $\partial_s\partial_i\eta$ in
\eqref{eq:finite-type-hypothesis} has $\delta_{si}=0$.
Substituting that equation for the derivatives of $v$ yields
\begin{equation}\label{eq:finite-type-system}
 \begin{aligned}
  \partial_s v_a&=\sum_{b=1}^n B_{as}^{\ b}v_b+\delta_{as}q,\\
  \partial_s q&=\sum_{b=1}^n C_s^{\ b}v_b+d_s q,
 \end{aligned}
\end{equation}
where the smooth coefficients are explicitly
\begin{equation}\label{eq:finite-type-coefficients}
 \begin{aligned}
 C_s^{\ b}
   &=\partial_i B_{si}^{\ b}-\partial_s B_{ii}^{\ b}
     +\sum_{a=1}^n
       \bigl(B_{si}^{\ a}B_{ai}^{\ b}-B_{ii}^{\ a}B_{as}^{\ b}\bigr),\\
 d_s&=B_{si}^{\ i}-B_{ii}^{\ s}.
 \end{aligned}
\end{equation}
This is the claimed homogeneous system for $v$ and $q$.

If $\eta=0$ on an open subset, then $v=0$ there, and a diagonal entry of
\eqref{eq:finite-type-hypothesis} also gives $q=0$ there.  Along any
piecewise smooth path in $U$, Equation~\eqref{eq:finite-type-system}
restricts to a homogeneous linear ordinary differential equation for
$(v,q)$.  Uniqueness with zero initial data propagates its vanishing along
the path.  An open connected subset of $\R^n$ is path connected, so
$v=q=0$ on $U$.  Finally, $\eta$ is constant on $U$, and its value on the
initial open subset is zero.
\end{proof}

\subsection{Completion of the coefficient recovery}

\begin{proof}[Proof of Theorem~\ref{thm:main}]
Suppose that the two displacement-to-traction maps agree, and use the
common exterior extensions furnished by Lemma~\ref{lem:common-extension}.
Propositions~\ref{prop:transport-matching} and
\ref{prop:comparison-rigidity} give equality of the transport matrices
on $\Ecal_\theta$ for every nonzero null $\theta$.  In particular,
\eqref{eq:transport-coefficients} gives
\[
 D_\theta\log(k_1/\ell_1)
   =D_\theta\log(k_2/\ell_2).
\]
Since null vectors span $\C^3$, the difference of these logarithms has
zero gradient.  Its exterior value is zero.  Hence $k_1/\ell_1=k_2/\ell_2$
on $\R^3$, and the functions
\begin{equation}\label{eq:common-ratios}
 c=\frac{m_1}{\ell_1}=\frac{m_2}{\ell_2}>0,
 \qquad b_* =\log\frac{k_1}{m_1}=\log\frac{k_2}{m_2}
\end{equation}
are well defined and smooth.  Indeed, $m_j/\ell_j=1-k_j/\ell_j$, and
$k_j/m_j=(k_j/\ell_j)/(1-k_j/\ell_j)$.

Equality of the bottom-left blocks on $\Ecal_\theta$ means
\[
 (Q_{1,\theta}-Q_{2,\theta})\cdot p=0
 \qquad(p\in H_\theta).
\]
The annihilator of $H_\theta=\theta^\perp$ for the nondegenerate
ambient bilinear form is $\C\theta$.  Since each $Q_{j,\theta}$ is
linear in $\theta$, there is a smooth ambient matrix $K(x)$ such that
\begin{equation}\label{eq:null-line-preservation}
 Q_{1,\theta}-Q_{2,\theta}=K(x)\theta\in\C\theta
 \qquad(\theta\cdot\theta=0).
\end{equation}
For fixed $x$, the proportionality factor in
\eqref{eq:null-line-preservation} is a holomorphic function on $\mathcal C$.
To see this at every point, choose a local holomorphic representative
$\theta$ and a component $\theta_a\ne0$; the factor is
$(K\theta)_a/\theta_a$.  These expressions agree on overlapping charts
and are unchanged by rescaling the representative.  Compactness of
$\mathcal C$ makes the factor constant in direction.  The spanning
property of null vectors then gives
\begin{equation}\label{eq:scalar-ambient-difference}
 K(x)=\kappa(x)\Id,\qquad \kappa(x)=\tfrac13\tr K(x).
\end{equation}
In particular, $\kappa$ is smooth.

Set
\[
 \eta=\log m_1-\log m_2,\qquad v=\nabla\eta=g_1-g_2,
 \qquad \beta=\nabla b_*,\qquad \alpha=\frac{c-1/2}{c},
\]
where $g_j=\nabla\log m_j$.  Because
$\nabla\log k_j=g_j+\beta$, the expression for $Q$ in
\eqref{eq:transport-coefficients} becomes
\[
 2Q_{j,\theta}
   =\bigl((c-1/2)g_jg_j^T+c g_j\beta^T
                 -c\nabla^2\log m_j\bigr)\theta.
\]
Subtract these identities, use
$g_1g_1^T-g_2g_2^T=g_1v^T+vg_2^T$, and apply
\eqref{eq:scalar-ambient-difference}.  Defining the smooth scalar
$q=-2\kappa/c$, we obtain
\[
 \nabla^2\eta
    =\alpha(g_1v^T+vg_2^T)+v\beta^T+q\Id.
\]
Equivalently,
\begin{equation}\label{eq:coefficient-hessian}
 \begin{aligned}
 \partial_i\partial_j\eta
   &=\sum_{a=1}^3 B_{ij}^{\ a}\partial_a\eta+\delta_{ij}q,\\
 B_{ij}^{\ a}
   &=\alpha(g_1)_i\delta_{ja}
        +\bigl(\alpha(g_2)_j+\beta_j\bigr)\delta_{ia}.
 \end{aligned}
\end{equation}
The coefficients $B_{ij}^{\ a}$ are smooth on $\R^3$.  They are fixed
functions determined by the two coefficient pairs under comparison.
The common exterior extension gives $\eta=0$ on a nonempty exterior open
set.  Lemma~\ref{lem:finite-type}, applied to
\eqref{eq:coefficient-hessian}, yields $\eta=0$ throughout $\R^3$.
Thus $m_1=m_2$.  Equation~\eqref{eq:common-ratios} then gives
$k_1=m_1e^{b_*}=m_2e^{b_*}=k_2$, and therefore
\[
 \mu_1=\mu_2,\qquad
 \lambda_1=k_1-m_1=k_2-m_2=\lambda_2.
\]
Restricting to $\Omega$ proves the theorem.
\end{proof}
